% Section 14.8, translated from sv/chapters/14/exercises.tex.

\section{Exercises}
\label{c14:sec:uppgifter}

\begin{aside}
\textbf{How to work with the exercises.} Parts~1 and~2 are done during the lesson: first on your
own, a few minutes per exercise, then together. Part~3 is extra exercises for further practice
after the lesson. Most of them can be done in your head or with pencil and paper.

\facitnotis{L14}
\end{aside}

% ------------------------------------------------------------------------------------------
\uppgiftsdel{Part 1}{Review exercises}

\uppgift{1.1}{Rules of differentiation}
Differentiate the following functions.
\begin{delar}(3)
  \task $f_1(x) = x^2 \cdot e^{4x}$
  \task $f_2(x) = x^3 \cdot \ln 2x$
  \task $f_3(x) = \dfrac{x^3}{\ln x^2}$
\end{delar}

% ------------------------------------------------------------------------------------------
\uppgiftsdel{Part 2}{New material}

\uppgift{2.1}{Energy stored in an inductor}
The current through an inductor depends on time and is given by the function
\[
  i(t) = 5t,
\]
where $t$ is the time in seconds. The current is measured in amperes (A). The inductance of the
inductor is $L = 0.20\,\text{H}$ (henry).

The power $p(t)$ developed in the inductor is given by the function
\[
  p(t) = L \cdot i(t) \cdot i'(t),
\]
where
\begin{itemize}
  \item $p(t)$ is the power developed in watts (W),
  \item $L$ is the inductance of the inductor in henries (H),
  \item $i(t)$ is the current through the inductor in amperes (A),
  \item $i'(t)$ is the rate of change of the current through the inductor in A/s.
\end{itemize}
The energy $w(t)$ stored in the inductor between time $0$ and time $t$ is given by the function
\[
  w(t) = \int_0^t p(t)\,dt
\]
and is measured in joules (J).
\begin{parts}
  \item Find $p(t)$.
  \item Find the antiderivative $P(t)$ of $p(t)$.
  \item Find an expression for the energy $w(t)$ in the inductor as a function of time.
  \item The inductor is uncharged at the start, that is, $w(0) = 0$. Find the constant of
    integration $C$.
  \item Find how much energy has been stored in the inductor during the first $3$ seconds
    (${0 \leq t \leq 3}$).
\end{parts}

% ------------------------------------------------------------------------------------------
\uppgiftsdel{Part 3}{Extra exercises}
The exercises below are extra practice, best done on your own after the lesson.

\uppgift{3.1}{Antiderivatives of polynomials}
Find the antiderivative.
\begin{delar}(3)
  \task $\displaystyle\int x^3\,dx$
  \task $\displaystyle\int 4x\,dx$
  \task $\displaystyle\int 5\,dx$
  \task $\displaystyle\int (2x + 3)\,dx$
  \task $\displaystyle\int (x^2 - 4x + 1)\,dx$
\end{delar}

\uppgift{3.2}{Antiderivatives of standard functions}
Find the antiderivative.
\begin{delar}(3)
  \task $\displaystyle\int e^x\,dx$
  \task $\displaystyle\int e^{3x}\,dx$
  \task $\displaystyle\int \frac{1}{x}\,dx$
  \task $\displaystyle\int \sin(x)\,dx$
  \task $\displaystyle\int \cos(2x)\,dx$
  \task $\displaystyle\int \sin(4x)\,dx$
\end{delar}

\uppgift{3.3}{Checking by differentiation}
\begin{parts}
  \item Show that $F(x) = \dfrac{x^3}{3}$ is an antiderivative of $f(x) = x^2$.
  \item Show that $F(x) = -\dfrac{\cos(2x)}{2}$ is an antiderivative of $f(x) = \sin(2x)$.
  \item Show that $F(t) = -2e^{-t/2}$ is an antiderivative of $f(t) = e^{-t/2}$.
  \item Why is the answer to an indefinite integral never unique?
\end{parts}

\uppgift{3.4}{Definite integrals: polynomials}
Calculate.
\begin{delar}(3)
  \task $\displaystyle\int_0^2 3x^2\,dx$
  \task $\displaystyle\int_1^3 2x\,dx$
  \task $\displaystyle\int_0^4 (x + 1)\,dx$
  \task $\displaystyle\int_{-1}^{1} x^2\,dx$
  \task $\displaystyle\int_0^3 (6 - 2x)\,dx$
\end{delar}

\uppgift{3.5}{Definite integrals: other functions}
Calculate.
\begin{delar}(3)
  \task $\displaystyle\int_0^1 e^x\,dx$
  \task $\displaystyle\int_0^{\pi} \sin(x)\,dx$
  \task $\displaystyle\int_0^{\pi/2} \cos(x)\,dx$
  \task $\displaystyle\int_1^{e} \frac{1}{x}\,dx$
  \task $\displaystyle\int_0^2 e^{-t}\,dt$
\end{delar}

\uppgift{3.6}{Area under a curve}
\begin{parts}
  \item Calculate the area under $f(x) = x$ for $0 \leq x \leq 4$.
  \item Check your answer to a) with the formula for the area of a triangle.
  \item Calculate the area under $f(x) = 3$ for $1 \leq x \leq 5$ and check it with a rectangle.
  \item Calculate the area under $f(x) = x^2$ for $0 \leq x \leq 2$.
\end{parts}

\uppgift{3.7}{Negative areas}
The function $f(x) = x - 2$ is given.
\begin{delar}(3)
  \task Calculate $\displaystyle\int_0^2 (x - 2)\,dx$.
  \task Calculate $\displaystyle\int_2^4 (x - 2)\,dx$.
  \task Calculate $\displaystyle\int_0^4 (x - 2)\,dx$.
  \task* Explain why the integral in c) is zero even though the curve does not lie on the
    $x$-axis.
\end{delar}

\uppgift{3.8}{The constant of integration from an initial condition}
It is known that $f'(x) = 4x - 3$ and that $f(0) = 5$.
\begin{parts}
  \item Find the general antiderivative.
  \item Find the constant of integration $C$.
  \item Write down $f(x)$.
  \item Calculate $f(2)$.
\end{parts}

\uppgift{3.9}{Voltage from its rate of change}
The voltage across a capacitor changes at the rate $u'(t) = 6t\,\text{V/s}$, and at the start
$u(0) = 2\,\text{V}$.
\begin{parts}
  \item Find the general antiderivative.
  \item Find the constant of integration.
  \item Write down $u(t)$.
  \item Calculate $u(4)$.
\end{parts}

\uppgift{3.10}{Charge from current}
The current through a conductor is given by $i(t) = 4t + 2$ amperes, and $q(0) = 0$.
\begin{parts}
  \item Find an expression for the charge $q(t)$.
  \item Calculate $q(3)$.
  \item Calculate $q(5)$.
  \item How much charge passes between $t = 3\,\text{s}$ and $t = 5\,\text{s}$?
  \item Check your answer to d) by calculating $\displaystyle\int_3^5 i(t)\,dt$.
\end{parts}

\uppgift{3.11}{Charge from a sinusoidal current}
The current is given by $i(t) = 2\sin(100\pi t)$ amperes.
\begin{parts}
  \item Find the antiderivative of $i(t)$.
  \item Calculate $\displaystyle\int_0^{0.01} i(t)\,dt$.
  \item Calculate $\displaystyle\int_0^{0.02} i(t)\,dt$.
  \item Explain the result in c).
\end{parts}

\uppgift{3.12}{Energy stored in an inductor}
The current through an inductor with $L = 0.4\,\text{H}$ is given by $i(t) = 3t$ amperes. The
power is given by the relationship $p(t) = L\,i(t)\,i'(t)$.
\begin{parts}
  \item Find $i'(t)$.
  \item Find $p(t)$.
  \item Find the energy $w(t) = \displaystyle\int_0^t p(\tau)\,d\tau$.
  \item Calculate the energy after $2$ seconds.
  \item Check your answer with the formula $w = \dfrac{Li^2}{2}$.
\end{parts}

\uppgift{3.13}{Energy dissipated in a resistor}
The power in a resistor is given by $p(t) = 5 + 2t$ watts for $0 \leq t \leq 10\,\text{s}$.
\begin{parts}
  \item Find the energy $w(t) = \displaystyle\int_0^t p(\tau)\,d\tau$.
  \item Calculate the total energy after $10$ seconds.
  \item How much energy is dissipated during the last $5$ seconds?
  \item Calculate the average power over the whole interval.
\end{parts}

\uppgift{3.14}{Mean value of a function}
The mean value of $f$ over the interval $[a, b]$ is given by
\[
  \frac{1}{b - a}\int_a^b f(x)\,dx.
\]
\begin{parts}
  \item Calculate the mean value of $f(x) = x^2$ over $[0, 3]$.
  \item Calculate the mean value of $f(x) = 4$ over $[1, 5]$.
  \item Calculate the mean value of $i(t) = 10\sin(100\pi t)$ over half a period,
    $[0, 0.01]$.
  \item What fraction of the amplitude is the mean value in c)?
\end{parts}

\uppgift{3.15}{Rules for integrals}
It is known that $\displaystyle\int_0^2 f(x)\,dx = 5$ and $\displaystyle\int_0^2 g(x)\,dx = 3$.
Calculate:
\begin{delar}(2)
  \task $\displaystyle\int_0^2 \bigl[f(x) + g(x)\bigr]\,dx$
  \task $\displaystyle\int_0^2 2f(x)\,dx$
  \task $\displaystyle\int_0^2 \bigl[3f(x) - 2g(x)\bigr]\,dx$
  \task $\displaystyle\int_2^0 f(x)\,dx$
\end{delar}

\uppgift{3.16}{Differentiation and integration as opposites}
\begin{parts}
  \item Differentiate $F(x) = x^3 - 2x$.
  \item Integrate the result from a).
  \item Why does a constant appear in b) that was not in $F(x)$?
  \item Calculate $\dfrac{d}{dx}\left[\displaystyle\int 5x^4\,dx\right]$.
\end{parts}

\uppgift{3.17}{Finding an unknown limit or constant}
\begin{parts}
  \item Find $a > 0$ such that $\displaystyle\int_0^a 2x\,dx = 25$.
  \item Find $a$ such that $\displaystyle\int_0^a 3\,dx = 12$.
  \item Find $k$ such that $\displaystyle\int_0^2 kx\,dx = 10$.
  \item A capacitor is charged with the constant current $i = 0.5\,\text{A}$. After how long has a
    charge of $2\,\text{C}$ passed?
\end{parts}

\uppgift{3.18}{Find the error}
Each line contains a common error. Explain the error and give the correct answer.
\begin{delar}(2)
  \task $\displaystyle\int x^2\,dx = 2x + C$
  \task $\displaystyle\int e^{2x}\,dx = 2e^{2x} + C$
  \task $\displaystyle\int \sin(x)\,dx = \cos(x) + C$
  \task $\displaystyle\int_1^2 3x^2\,dx = \Bigl[x^3\Bigr]_1^2 = 8$
  \task* A definite integral needs a constant of integration.
  \task* $\displaystyle\int \frac{1}{x}\,dx = \ln x + C$ holds for all $x \neq 0$.
\end{delar}
